% J. L. Doob, "Probability in function space", Bull. Amer. Math. Soc. 53 (1947), first page, reset in Mills 8A.
\documentclass[11pt]{article}
\usepackage{bulletin1947}
\begin{document}
\setcounter{page}{15}
\received{An address delivered at the invitation of the Committee to Select Hour Speakers for Annual and Summer Meetings before the Summer Meeting of the Society on August 22, 1946; received by the editors August 23, 1946.}
\articletitle{PROBABILITY IN FUNCTION SPACE}{j. l. doob}
\head{Introduction.} The mathematical theory of probability is now ordinarily formulated in terms of measure theory. The formalization of a succession of $n$ trials is in terms of $n$-dimensional measure; that of a sequence of trials is in terms of (denumerably) infinite-dimensional measure; that of a continuous set of trials (a set dependent on a continuous time parameter say) is in terms of measure in function space. The latter case is the subject of this lecture.

There still is a residual tendency among mathematicians to deal with distribution functions rather than with the chance variables which have those distribution functions; in this way probability theorems are reduced to theorems on monotone functions, convolutions, and so on. This preference was once justified by a natural distrust of the mathematical foundations of probability theory, so that any discussion of chance variables seemed to be nonmathematical, physics, or at best mathematical physics rather than pure mathematics, and therefore not rigorous. By now this distrust of chance variables is no longer justified; in practice, although it has almost disappeared when finite or denumerably infinite sequences of chance variables are studied, its presence is still noticeable in studies of continuous families of chance variables. Thus there is now no hesitancy in evaluating the probability that a series whose terms depend on chance will converge, but there is still considerable hesitancy in discussing the probability that a function whose values depend on chance should have a limiting value at $\infty$, be continuous, or be Lebesgue integrable. It is the purpose of the present lecture to dispel some inhibitions by summarizing the basic knowledge on measure theory in function space, as applied to probability. In the course of doing this, some problems still to be solved will be discussed. Although the theory is mathematically rigorous, it cannot yet be considered fully satisfactory. Even mathematicians, although they can live on rigor alone, do not invariably enjoy it.

The general point of view adopted in analyzing a succession of trials, depending on the parameter $t$ (usually representing time), is the following. A set $T$ of values of $t$ is preassigned, the times when events occur, or measurements are made. This set of times may be finite or infinite. At each such moment a number $x(t)$ is obtained.
\end{document}
